\subsection{本质可积函数}

设 F 是一个实 Banach 空间；回忆一下，空间 \(F_{F}^{p}\) （第四章，§3 第 3 号）与 \(L_{F}^{p}\) （第四章，§3 第 4 号定义 2）的元素都是 \(\mu\) -适度函数（第四章，§5 第 6 号引理 1）；仍用 \(N_{F}\) 表示 T 到 F 中的可忽略映射所成的空间，我们将引进 T 到 F 中的局部可忽略映射所成的空间 \(N_{F}^{\infty}\) 。

引理. --- 设 g 与 \(g'\) 是两个取值于 F 的 \(\mu\) -适度映射；若 g 与 \(g'\) 局部几乎处处等于同一个函数 f，则 \(g = g'\) 几乎处处成立。

设 D 是由 \(t \in T\) （使 \(\mathbf{g}(t) \neq \mathbf{g}'(t)\) 者）组成的集；D 局部可忽略且适度，因此可忽略（命题 7 推论 1）。

我们将用 \(\overline{\mathcal{F}}_F^p (\mathrm{T},\mu)\) （若不致混淆，也简记作 \(\overline{\mathcal{F}}_{\mathrm{F}}^{p}(\mu),\overline{\mathcal{F}}_{\mathrm{F}}^{p}\) ）表示映射 \(\mathbf{f}\) （\(\mathrm{T}\) 到 \(\mathrm{F}\) 中的映射）所成之集，其中存在函数 \(\mathbf{g}\in \mathcal{F}_{\mathrm{F}}^{p}\) 局部几乎处处等于 \(\mathbf{f}\) 。由引理，数 \(\mathrm{N}_p(\mathbf{g})\) 只依赖于 \(\mathbf{f}\) ，故我们将写 \(\overline{\mathrm{N}}_p(\mathbf{f}) = \mathrm{N}_p(\mathbf{g})\) 。函数 \(\overline{\mathrm{N}}_p\) 显然是 \(\overline{\mathcal{F}}_F^p\) 上的一个半范数，且我们总假定 \(\overline{\mathcal{F}}_F^p\) 赋有由 \(\overline{\mathrm{N}}_p\) 定义的拓扑。0 在此拓扑下的闭包是空间 \(\mathcal{N}_{\mathrm{F}}^{\infty}\) ；关系式 \(\overline{\mathcal{F}}_{\mathrm{F}}^{p} = \mathcal{F}_{\mathrm{F}}^{p} + \mathcal{N}_{\mathrm{F}}^{\infty}\) 与 \(\mathcal{N}_{\mathrm{F}}^{\infty}\cap \mathcal{F}_{\mathrm{F}}^{p} = \mathcal{N}_{\mathrm{F}}\) （引理）表明，赋范空间 \(\overline{\mathcal{F}}_{\mathrm{F}}^{p} / \mathcal{N}_{\mathrm{F}}^{\infty}\) 可典范地等同于 \(\mathcal{F}_{\mathrm{F}}^{p} / \mathcal{N}_{\mathrm{F}}\) ，后者是完备的（第四章，§3 第 3 号命题 5）；因此 \(\overline{\mathcal{F}}_F^p\) 本身是完备的。

类似地，我们将用 \(\overline{\mathcal{L}}_F^p (\mathrm{T},\mu)\) （或 \(\overline{\mathcal{L}}_F^p (\mu)\) ，或 \(\overline{\mathcal{L}}_F^p\) ）表示子空间 \(\mathcal{L}_F^p +\mathcal{N}_F^\infty\) （\(\overline{\mathcal{F}}_F^p\) 的） ；也可以把 \(\overline{\mathcal{L}}_F^p\) 刻画为 \(\overline{\mathcal{F}}_F^p\) 中由可测映射组成的子空间（第四章，§5 第 6 号定理 5）。

赋范空间 \(\overline{L}_{F}^{p}/N_{F}^{\infty}\) 可典范地等同于 \(L_{F}^{p}\) ；因此 \(\overline{L}_{F}^{p}\) 是完备的。其元素称为 p 次幂本质可积函数，这一术语的合理性由下述命题得到证明：

命题 9. --- 为使 T 到 F 中的映射 f 属于 \(\overline{\mathcal{F}}_F^p\) （相应地，属于 \(\overline{\mathcal{L}}_F^p\) ），必须且只须（相应地，只须 f 可测且）

\[
\mu^ {\bullet} (| \mathbf {f} | ^ {p}) <   + \infty .
\]

于是有 \(\overline{\mathrm{N}}_p(\mathbf{f}) = (\mu^\bullet (|\mathbf{f}|^p))^{1 / p}\) 。

显然我们可以只限于关于 \(\overline{F}_{F}^{p}\) 的断言。若 f 属于 \(\overline{F}_{F}^{p}\) ，设 g 是属于 \(F_{F}^{p}\) 且局部几乎处处等于 f 的函数；则 \(|f|^{p}=|g|^{p}\) 局部几乎处处成立，因此

\[
\mu^ {\bullet} (| \mathbf {f} | ^ {p}) = \mu^ {\bullet} (| \mathbf {g} | ^ {p}) = \mu^ {*} (| \mathbf {g} | ^ {p}) <   + \infty
\]

（命题 1 a) 与命题 7），且另一方面，由 \(\overline{\mathbf{N}}_p\) 的定义，

\[
\overline {{\mathrm{N}}} _ {p} (\mathbf {f}) = \mathrm{N} _ {p} (\mathbf {g}) = \left(\mu^ {*} (| \mathbf {g} | ^ {p})\right) ^ {1 / p}.
\]

反之，设 \(\mu^{\bullet}(|\mathbf{f}|^{p}) < +\infty\) ；则存在适度集 A，使 f 在 T-A 中局部几乎处处为零（命题 7）。函数 \(f\varphi_{A}\) 局部几乎处处等于 f，且使 \(\mathrm{N}_{p}(\mathbf{f}\varphi_{\mathrm{A}}) = \overline{\mathrm{N}}_{p}(\mathbf{f}) < +\infty\) ，因此它属于 \(F_{F}^{p}\) ，从而 \(f \in F_{F}^{p}\) 。

推论. --- 为使 f 属于 \(L_{F}^{p}\) ，必须且只须 f 属于 \(\overline{L}_{F}^{p}\) 且是适度的。

定义 3. -- \(\overline{\mathcal{L}}_F^1\) 的元素称为取值于 F 的本质 \(\mu\) -可积函数。把映射 \(\widetilde{\mathbf{f}} \mapsto \mu(\mathbf{f})\) （\(L sub F superscript 1\) 到 F 中的）与 \(\overline{\mathcal{L}}_F^1\) 到 \(L sub F superscript 1\) 上的典范映射复合，就得到 \(\overline{\mathcal{L}}_F^1\) 到 F 中的一个连续线性映射，它是映射 \(\mathbf{f} \mapsto \int \mathbf{f} d\mu\) （\(\mathcal{L}_F^1\) 到 F 中的映射）的延拓。此映射的值仍记作 \(\int \mathbf{f} d\mu\) 或 \(\mu(\mathbf{f})\) （对 \(\mathbf{f} \in \overline{\mathcal{L}}_F^1\) ），并称这个元素为 \(\mathbf{f}\) 关于 \(\mu\) 的积分。

两个局部几乎处处相等的本质可积函数有相同的积分。对每个有限且本质可积的函数 \(f \geqslant 0\) ，有 \(\int^{\bullet} f d\mu = \int f d\mu\) 。若 A 是一个其特征函数本质可积的集，则称 A 为本质 \(\mu\) -可积集；\(\int \varphi_{A} d\mu\) 也记作 \(\mu(A)\) ，并仍称为 A 的测度。

若函数 \(\mathbf{f}\) （取值于 \(\mathbf{F}\) ）只在 \(\mathbf{T}\) 中局部几乎处处有定义，我们仍称 \(\mathbf{f}\) 是本质可积的，如果它局部几乎处处等于一个处处有定义且可积的函数 \(\mathbf{f}_1\) ；这时令

\[
\int \mathbf {f} d \mu = \int \mathbf {f} _ {1} d \mu ,
\]

且此定义不依赖于处处有定义且局部几乎处处等于 f 的可积函数 \(f_{1}\) （引理）。对取值于 \(\overline{R}\) 且局部几乎处处有定义并有限的函数，类似地定义本质可积函数的概念。

读者不难把第四章 §4 中关于可积函数的结果推广到本质可积函数，只需把命题叙述中的“几乎处处”换成“局部几乎处处”。例如我们指出不等式

\[
\left| \int \mathbf {f} d \mu \right| \leqslant \int | \mathbf {f} | d \mu ,\tag{3}
\]

它对每个取值于 Banach 空间的本质可积函数 f 成立。

命题 10. --- 设 \(\mathfrak{K}\) 是 T 的紧子集的一个 \(\mu\) -稠密集。a) 若 \(f\) 是 \(\geqslant 0\) 的数值函数，则

\[
\mu^ {\bullet} (f) = \sup _ {K \in \mathfrak {K}} \mu^ {*} (f \varphi_ {K}).\tag{4}
\]

\begin{enumerate}
\def\labelenumi{\alph{enumi})}
\setcounter{enumi}{1}
\tightlist
\item
  若 \(\mathbf{f}\) 是取值于 Banach 空间 \(\mathbf{F}\) 的本质可积函数，则
\end{enumerate}

\[
\int \mathbf {f} d \mu = \lim _ {\mathfrak {K}} \int \mathbf {f} \varphi_ {\mathrm{K}} d \mu ,
\]

极限是沿对 \(\subset\) 有向的集 \(\mathfrak{K}\) 取的。

为建立 a)，只需证明：对 T 的每个紧子集 L，\(\int^{*}f\varphi_{L}d\mu=\sup_{K}\int^{*}f\varphi_{K}d\mu\) ，其中 K 取遍属于 K 的 L 的子集所成之集。由于 L 是一个可忽略集与 K 中元素的一个递增序列 \((K_{n})\) 之并（第四章，§5 第 8 号命题 12），这由上积分中的极限过渡定理（第四章，§1 第 3 号定理 3）得出。

现在设 f 属于 \(\overline{L}_{F}^{1}\) ；设 \(\varepsilon\) 是一个数 \textgreater0，并设 K 是 K 的一个元素，使

\[
\int | \mathbf {f} | \varphi_ {\mathrm{K}} d \mu \geqslant \int | \mathbf {f} | d \mu - \varepsilon
\]

（这样的 K 存在，由 a)）。于是，对包含 K 的每个紧集 H，

\[
\left| \int \mathbf {f} d \mu - \int \mathbf {f} \varphi_ {\mathrm{H}} d \mu \right| \leqslant \int | \mathbf {f} | \varphi_ {\mathbf {C H}} d \mu \leqslant \int | \mathbf {f} | \varphi_ {\mathbf {C K}} d \mu \leqslant \varepsilon .
\]

复 Banach 空间与复测度情形的推广。设 F 是一个复 Banach 空间；滥用记号，F 的底实 Banach 空间也记作 F。Banach 空间 \(\overline{\mathcal{L}}_{\mathrm{F}}^{p}(\mathrm{T},\mu)\) 这时可赋以一个自然的复 Banach 空间结构，且必须明确指出使用的是此空间的实结构还是复结构。在本章中，除非明确说明相反情形，总理解为实结构。

设 \(\theta\) 是一个复测度；我们令 \(\overline{\mathcal{L}}_F^p (\mathrm{T},\theta) = \overline{\mathcal{L}}_F^p (\mathrm{T},|\theta |)\) ；若 \(\mathbf{F}\) 是复 Banach 空间，可以作出与上面相同的注记。特别地，函数 \(\mathbf{f}\) （取值于 \(\mathbf{F}\) ）将称为对 \(\theta\) 本质可积的，如果它对 \(|\theta |\) 本质可积。于是命题 10 的断言 b) 立即推广到复测度。

